\section{Introduction}
\label{sec:introduction}

For a discrete group $\Gamma$, its integral cohomological dimension $\cdim\Gamma$ is the projective dimension of the trivial left $\Z[\Gamma]$-module $\Z$: it measures the length of the shortest projective resolution of $\Z$. Its geometric dimension $\gdim\Gamma$ is the least dimension of a connected CW complex with fundamental group $\Gamma$ and contractible universal cover. Such a complex is a \emph{classifying space}, or a $K(\Gamma,1)$. These are ordinary dimensions; no family of subgroups or proper-action variant is involved. The cellular chains of a contractible universal cover give a free resolution, so $\cdim\Gamma\le\gdim\Gamma$.

The Eilenberg--Ganea theorem of 1957 identifies these dimensions when $\cdim\Gamma\ge3$. In dimension one, Stallings proved that a finitely generated group of integral cohomological dimension one is free, and Swan removed the finite-generation hypothesis \cite{Stallings1968,Swan1969}; a free group has a one-dimensional classifying space. The remaining dimension-two conjecture asks whether every group of integral cohomological dimension two has a two-dimensional classifying space \cite{EG1957}. We answer this question negatively with a finitely generated residually finite group. Here \emph{residually finite} means that every nonidentity element remains nonidentity in some finite quotient.

Our group comes from a finite simplicial complex. A simplicial complex is \emph{flag} if every finite set of pairwise adjacent vertices spans a simplex. For such a complex $L$ with vertex set $V$, the associated \emph{right-angled Artin group} is
\[
 A_L=\left\langle a_q\ (q\in V)\ \middle|\
 [a_q,a_r]=1\text{ if }\{q,r\}\text{ is an edge of }L\right\rangle.
\]
The \emph{height homomorphism} $\lambda:A_L\to\Z$ sends every standard generator $a_q$ to $1$. Its kernel is a Bestvina--Brady group. Acyclicity below always means vanishing reduced integral homology.

\begin{theorem}
\label{thm:main}
There is a finitely generated residually finite group $G$ such that
\[
 \cdim G=2,\qquad \gdim G=3.
\]
One may take $G$ to be the height kernel in the right-angled Artin group associated to a finite acyclic flag triangulation of the presentation complex
\begin{equation}
\label{eq:seed-presentation-intro}
 \langle x,y\mid x^2=y^5,\ x^2=(xy^{-1})^3\rangle.
\end{equation}
No two-dimensional classifying space for this $G$ exists, regardless of the number of its cells.
\end{theorem}

Section~\ref{sec:geometry} constructs the triangulation and proves the stated algebraic and geometric properties of its height kernel. Finite generation also makes $G$ countable, so restricting the conjecture to countable groups does not restore it.

Proposition~\ref{prop:model} also supplies a length-two resolution of the trivial $\Z[G]$-module $\Z$ by finitely generated free $\Z[G]$-modules. A group admitting such a finite-length free resolution is said to be of \emph{type $\mathrm{FL}$}. Consequently, $G$ is of \emph{type $\mathrm{FP}_\infty$}: it admits a projective resolution with finitely generated terms in every degree. Here the length-two free resolution extends to such a projective resolution by taking all higher terms to be zero. The group is not finitely presented: $L$ has nontrivial fundamental group by Lemma~\ref{lem:geometric-seed}, and the Bestvina--Brady finite-presentability criterion applies \cite[Main Theorem, part~(3)]{BB1997}.

\subsection*{History and the obstruction}

Bestvina and Brady used height kernels to distinguish homological from homotopical finiteness conditions \cite[Main Theorem]{BB1997}. Their acyclic two-dimensional examples provide candidates for the dimension-two problem \cite[Example~6.3(3)]{BB1997}. For suitable flag triangulations of a spine of the Poincar\'e homology sphere, they proved that at least one of the Eilenberg--Ganea and Whitehead conjectures fails \cite[Theorem~8.7]{BB1997}. Whitehead's asphericity question asks whether every connected subcomplex of an aspherical two-dimensional CW complex is aspherical \cite{Whitehead1941}; see also \cite[Section~8]{BB1997}. The classical alternative does not itself determine the geometric dimension of a height kernel. Leary and Petrosyan give a recent conditional comparison between Coxeter and Bestvina--Brady candidates for the ordinary problem \cite[Section~5.6, Conjecture~5.20 and Corollary~5.21]{LP2026}.

Dicks and Leary describe height kernels by edge generators and power relators associated to cycles \cite[Theorem~1 and Proposition~2]{DicksLeary1999}. Howie studied related Bestvina--Brady examples through plus constructions and obtained finite-rank free relation modules \cite{Howie1999}. These works provide the construction and relation-module setting for the present paper. The interface needed here is the particular free basis represented by triangular boundary chains; its exact form is displayed below and proved directly in Proposition~\ref{prop:model}.

The obstacle is that a free action on an acyclic two-complex supplies a short free resolution, whereas a classifying space requires a contractible universal cover. An acyclic complex need not be contractible. Proving that one acyclic model is noncontractible therefore leaves open the possibility of a different two-dimensional classifying space. We compare the cycle module of the acyclic model with the presentation arising from \emph{any} hypothetical two-dimensional classifying space. The comparison permits arbitrary sets of generators and relators in that hypothetical model.

\subsection*{The geometric model and the proof interfaces}

For the complex $L$ in Theorem~\ref{thm:main}, form the union of the coordinate subtori in $(S^1)^V$ indexed by the simplices of $L$, using the common basepoint in all other coordinates. This is a cubical model for $A_L$ with a contractible three-dimensional universal cover $E$. Identify its vertices with $A_L$. The height homomorphism extends affinely over its cubes, and
\[
 X=\lambda^{-1}(0)\subset E
\]
is a two-dimensional level complex with vertex set $G$. The group $G=\ker\lambda$ acts freely on $X$ by left translation, with one vertex orbit and finitely many cell orbits. Proposition~\ref{prop:model} proves that $X$ is acyclic and that $G$ is finitely generated and residually finite, with $\cdim G=2$ and $\gdim G\le3$. These are the geometric inputs to the obstruction; their proofs are included in Section~\ref{sec:geometry}.

Here is the exact algebraic information supplied by $X$. Choose the lift based at $1$ of one oriented edge in each $G$-orbit. Its endpoint is the value of the corresponding letter in a finite generating alphabet $S$ for $G$; distinct edge orbits retain distinct letters. Reading the boundaries of representatives of the triangular two-cell orbits gives a finite indexed list $\mathcal D$ of words on $S$. Write parenthesized superscripts for direct sums. For $\Lambda=\Z[G]$, let $\partial_1:\Lambda^{(S)}\to\Lambda$ be the cellular edge boundary and let $[z]$ be the edge chain of a word $z\in\mathcal D$. The canonical map
\[
 \Lambda^{(\mathcal D)}\longrightarrow\ker\partial_1,
 \qquad e_z\longmapsto[z],
\]
is an isomorphism. Thus the specified boundary chains form a free basis, with both spanning and independence. The words in $\mathcal D$ are trivial in $G$, but they need not be a presentation of $G$; this distinction is the reason for the comparison argument.

We use labels in $H=\SU(2)$, identified with the unit quaternions. An edge labeling assigns an element of $H$ to each oriented edge and the inverse element to its reverse. A \emph{path product} multiplies the labels in traversal order; for a closed path we also call this its \emph{holonomy}. A labeling is \emph{flat on the triangles} when every triangular boundary has product $1$. The two main interfaces give incompatible conclusions about these labels:
\begin{enumerate}
\item Theorem~\ref{thm:presentation} shows that a hypothetical two-dimensional $K(G,1)$ would force every sufficiently small labeling, flat on all translates of the words in $\mathcal D$, to have trivial product along every closed path.
\item Proposition~\ref{prop:small-labels} constructs, for every $\epsilon>0$, labels of distance less than $\epsilon$ from $1$ that are flat on every triangle of $X$ and have nontrivial product along some closed path.
\end{enumerate}
Choosing the second labeling below the threshold in the first statement excludes every two-dimensional classifying space. The quotient $E/G$ supplies a three-dimensional one.

\subsection*{How the comparison and transport work}

Boundary chains record signed edge traversals but forget their order, so equal chains need not give equal path products in $H$. Theorem~\ref{thm:presentation} overcomes this loss of information using an actual aspherical presentation. After adding generators with defining relators, we change the relator words so that a finite distinguished list has exactly the boundary chains $[z]$ supplied by $X$; all additional generators represent the identity of $G$ and therefore give loop edges. The word identities in this presentation yield exact translated signed counts. Once the remaining relator equations hold, these counts cancel the linear discrepancy between the distinguished products and the triangular products. Commutators of elements close to $1$ contribute only a quadratic error. The construction works even if the hypothetical presentation has infinitely or uncountably many generators and relators.

This estimate prevents the distinguished products from reaching a fixed positive distance from $1$ while the prescribed edge labels remain sufficiently small. On a finite quotient, each finite collection of the remaining relator equations has independent vectors recording the total exponents of the unknown labels. The word-equation degree method of Gerstenhaber and Rothaus \cite{GR1962}, with related representation-theoretic applications in Fr\o yshov \cite{Froyshov2013}, supplies the starting point for solving these systems. At the coefficient-free starting system, Lemma~\ref{lem:localized-degree} gives nonzero degree on the conjugation-invariant region where the distinguished products are small. The quadratic barrier preserves this degree as the prescribed labels move to their required values.

Residual finiteness allows us to copy any finite collection of prescribed labels to a suitable quotient. Compactness then gives labels on all additional edges of the original graph. Since the original triangular products are exactly $1$, the same quadratic estimate forces the distinguished products to be exactly $1$ as well. All relators of an actual presentation now vanish, so every closed path has trivial product. The original labeling need not descend to any single finite quotient.

The transport construction contradicts this conclusion. Reduced words in $A_L$ give maps from the graph of $X$ to $L$ whose edge and triangle-boundary images are uniformly small, while special long paths map exactly around prescribed loops of $L$. The seed in Equation~\eqref{eq:seed-presentation-intro} supplies a nontrivial representation $\pi_1(L)\to\SU(2)$. Pulling its transport back along these maps gives small edge labels: the small triangle images force trivial triangular products, and an exactly traced loop detected by the representation retains nontrivial holonomy. Section~\ref{sec:transport} constructs the maps and labels explicitly.

\subsection*{Conventions and organization}

We measure distance in $H$ by the quaternion norm $\abs{a-b}$; this metric is bi-invariant. Our traversal-order convention makes the product along a path followed by another path the product of their respective labels in that order. We use
\begin{equation}
\label{eq:quaternion-basic}
 \abs{a_1\cdots a_k-1}\le\sum_{j=1}^k\abs{a_j-1},
 \qquad
 \abs{ab-ba}\le2\abs{a-1}\abs{b-1}.
\end{equation}
The first inequality follows by telescoping; the second follows by expanding
$ab-ba=(a-1)(b-1)-(b-1)(a-1)$.

Group rings act from the left. Generator edges run from $g$ to $gs$, and translating a path by $g$ means left translation. We work in ordinary set theory with choice, including compactness of products indexed by arbitrary sets.

Section~\ref{sec:degree} proves the localized degree statement. Section~\ref{sec:presentation} derives the presentation obstruction, including its arbitrary-cardinality comparison and compactness steps. Section~\ref{sec:geometry} constructs the seed, the group and the acyclic level model, and Section~\ref{sec:transport} constructs the small labels and completes the contradiction. Standard CW topology, covering-space theory, and finite-dimensional degree theory enter with their stated hypotheses at the relevant steps.
